\subsection{Dual Koszul cocycles and the prime-pair condition}
\label{depth:dual-cocycles}

The depth bound can also be approached by constructing classes in the
dual Koszul complex. At a variable cutoff write $R_m=R^{<m}$ and
$G_m=G^{<m}$, and let $\mathbb K_{\kk}^{<m}$ be the Koszul resolution of
$\kk$ over $R_m$. Put
\[
 C_m=\operatorname{Hom}_{R_m}(\mathbb K_{\kk}^{<m},G_m).
\]
Write $C_{m,n,d}$ for its complex in internal bidegree $(n,d)$ and
$C^i_{m,n,d}$ for the term in cohomological degree $i$.
In this subsection we use $\varepsilon_q$ for the dual exterior
generators, of cohomological degree one and internal bidegree
$(-q,-1)$. With the transpose convention for the codifferential,
\[
 \mathrm{d}^\star\omega=\theta_m\wedge\omega,
 \qquad \theta_m=\sum_{q\in\A_m}x_q\varepsilon_q.
\]
Thus $\mathrm{d}^\star$ preserves internal bidegree, and
$\mathrm H^i(C_m)=\operatorname{Ext}_{R_m}^i(\kk,G_m)$.
The least $i$ with nonzero cohomology is $\delta^{<m}$.
If $s_m=\sum_{q\in\A_m}q$, finite Koszul self-duality gives
\begin{equation}\label{depth:dual-betti}
 \dim_\kk\mathrm H^i(C_m)_{n,d}
 =\beta^{<m}_{|\A_m|-i,\,s_m+n,\,|\A_m|+d}.
\end{equation}

Suppose $m\ge11$ and every prime $5\le q\le m$ has a representation
$q+1=p_1+p_2$ by odd primes. Both summands are at most $q-2$, so they
belong to $\A_m$. Choose $p_1\le p_2$ with $p_1p_2$ maximal, using
the product order of Proposition~\ref{moves:order}, and set
\[
 y_q=x_{p_1}x_{p_2}\quad\text{in }G_m.
\]
The maximum-product choice fixes a representative; the class $y_q$
is independent of this choice by the defining quadratic relations.
In particular,
\[
 y_5=x_3^2,\qquad y_7=x_3x_5,\qquad y_{11}=x_5x_7.
\]
Consider the cochain
\begin{equation}\label{depth:dual-family}
 \omega_m=\varepsilon_3\wedge
       \sum_{\substack{5\le q\le m\\q\text{ prime}}}y_q\varepsilon_q
       \in C^2_{m,-2,0}.
\end{equation}

\begin{proposition}\label{depth:dual-family-criterion}
The cochain $\omega_m$ is closed if and only if
\begin{equation}\label{depth:dual-cubic-identities}
 x_qy_{q'}=x_{q'}y_q\quad\text{in }G_m
 \qquad(5\le q<q'\le m).
\end{equation}
Equivalently, if $q+1=p_1+p_2$ and $q'+1=p'_1+p'_2$ are the
chosen representations, then
\[
 x_q(x_{p'_1}x_{p'_2})=x_{q'}(x_{p_1}x_{p_2})\quad\text{in }G_m.
\]
Whenever it is closed, its class is nonzero over every field. In
particular,
\[
 \beta^{<m}_{|\A_m|-2,\,s_m-2,\,|\A_m|}\ge1,
 \qquad \delta^{<m}\le2.
\]
\end{proposition}
\begin{proof}
Exterior multiplication gives
\[
 \mathrm{d}^\star\omega_m=
 \sum_{5\le q<q'\le m}(x_{q'}y_q-x_qy_{q'})
        \varepsilon_3\varepsilon_q\varepsilon_{q'},
\]
which proves the criterion. A possible primitive in internal
bidegree $(-2,0)$ has the form $\omega'=\sum_q f_q\varepsilon_q$, with
$f_q\in(G_m)_{q-2,1}$. Since $1$ and $9$ are not prime labels,
$f_3=f_{11}=0$. The coefficient of
$\varepsilon_3\varepsilon_{11}$ in every such $\mathrm{d}^\star\omega'$
is therefore $x_3f_{11}-x_{11}f_3=0$. In $\omega_m$ that coefficient is the nonzero
monomial class $x_5x_7$. The component basis of
Proposition~\ref{setup:components} proves its nonvanishing over every
field. Equation~\eqref{depth:dual-betti} and the depth characterization
give the final assertions.
\end{proof}

\begin{example}[Finite cocycles]\label{depth:dual-finite}
At cutoff $12$ the cocycle is
\[
 \omega_{12}=x_3^2\varepsilon_3\varepsilon_5
       +x_3x_5\varepsilon_3\varepsilon_7
       +x_5x_7\varepsilon_3\varepsilon_{11}.
\]
Its closure follows from $x_3x_7=x_5^2$ and $x_3x_{11}=x_7^2$.
For example,
$x_5y_{11}=x_5^2x_7=x_3x_7^2=x_3^2x_{11}=x_{11}y_5$.
The following choices extend the displayed expression by one term
whenever a new prime is adjoined:
\[
\begin{array}{c|ccccc}
 q&13&17&19&23&29\\ \hline
 y_q&x_7^2&x_7x_{11}&x_7x_{13}&x_{11}x_{13}&x_{13}x_{17}
\end{array}
\]
Thus $\omega_{14}=\omega_{12}+x_7^2\varepsilon_3\varepsilon_{13}$ and
$\omega_{18}=\omega_{14}+x_7x_{11}\varepsilon_3\varepsilon_{17}$, and
similarly at cutoffs $20$, $24$, and $30$.

The supplement verifies~\eqref{depth:dual-cubic-identities} at every
odd-prime prefix through $101$. Hence $\omega_m$ is a nonboundary cocycle
over every field for every integer $11\le m\le102$.
The verification enumerates cubic monomials and their binary-move
components at each support; it does not replace $G_m$ by its toric
quotient. The coefficient test in the preceding proof certifies
nontriviality independently of the field.

At the smaller cutoff $8$ there is instead a linear-coefficient class
\[
 \omega_8=x_3\varepsilon_3\varepsilon_5+x_5\varepsilon_3\varepsilon_7
       \in C^2_{8,-5,-1}.
\]
Its codifferential is
$\mathrm{d}^\star\omega_8=(x_3x_7-x_5^2)\varepsilon_3\varepsilon_5\varepsilon_7=0$.
The only possible primitive in this bidegree is a scalar multiple
of $\varepsilon_5$; the codifferential of every nonzero such multiple
has a nonzero $\varepsilon_5\varepsilon_7$ coefficient, so the displayed class is
nonzero.
\end{example}

\begin{conjecture}[Infinite dual Koszul cocycle]\label{depth:dual-conjecture}
For every odd prime $q\ge5$, there are odd primes $p_1,p_2$ with
$q+1=p_1+p_2$. Let $y_q=x_{p_1}x_{p_2}\in G_{q+1,2}$ be the
resulting quadratic class in the full binary Goldbach ring. Then
\[
 x_qy_{q'}=x_{q'}y_q\quad\text{in }G\qquad(5\le q<q').
\]
\end{conjecture}

For infinitely many variables,
$C^i=\operatorname{Hom}_R(\mathbb K_{\kk,i},G)$ is a product over
the $i$-element exterior supports. Each coordinate of $\mathrm{d}^\star$ uses
only finitely many input coordinates. Under the conjecture,
\[
 \omega=\varepsilon_3\wedge\sum_{q\ge5}y_q\varepsilon_q
\]
therefore defines a cocycle in $C^2$. It also belongs to the
completion of the direct-sum cochain module in the Goldbach weight
of its coefficients: the coefficient $y_q$ has weight $q+1$.
This is a coefficient filtration; every whole term has dual internal
bidegree $(-2,0)$. The same $\varepsilon_3\varepsilon_{11}$
coefficient test proves that $\omega$ is not a boundary, even in this
completion. In particular, the conjecture gives
$\operatorname{Ext}_R^2(\kk,G)_{-2,0}\ne0$.

The conjecture concerns identities in the full ring. Requiring every
$\omega_m$ to be closed already in $G^{<m}$ is stronger; the finite
examples have this additional property. A cubic identity of weight
$q+q'+1$ may in general use primes as large as $q+q'-5$ along a move
path, beyond the original labels $q,q'$.

\begin{proposition}\label{depth:dual-prime-pair}
Conjecture~\ref{depth:dual-conjecture} implies that every prime
$q\ge7$ has a representation
\[
 q+3=p_1+p_2,\qquad p_1,p_2>3\text{ prime}.
\]
Consequently it implies $I^{<m}:x_3^\infty=J^{<m}$ at every
variable cutoff $m>5$.
\end{proposition}
\begin{proof}
The identity with the label $5$ is $x_5y_q=x_3^2x_q$.
If $q+3$ has no representation by primes greater than $3$, the
monomial $x_3^2x_q$ is isolated under binary moves: its only pair
sums are $6$ and $q+3$, each uniquely represented. For $q\ge7$ the
monomial $x_5y_q$ is different, since it contains $x_5$.
The component basis therefore contradicts the required equality.
Moreover, every prime occurring along a cubic path of weight $q+6$
is at most $q$, so larger cutoffs cannot repair a failure of this
particular identity. The saturation conclusion follows from the
prime-pair criterion in Appendix~\ref{depth:prime-variable}.
\end{proof}

The relation-group theorem does not supply the cubic identities in
Conjecture~\ref{depth:dual-conjecture}. When $y_q,y_{q'}$ exist,
$x_qy_{q'}-x_{q'}y_q$ belongs to $J$, so that theorem supplies a common
monomial multiplier annihilating the difference in $G$.
It does not permit cancellation of that multiplier. Even saturation
by powers of $x_3$ only gives annihilation by some power, rather
than equality in $G$. Thus the conjecture asks for specific cubic
binary-move paths beyond the stable connectivity already proved.
